\section{Results (Preliminary / Forthcoming)}
\label{sec:results}

\note{This section presents the planned result schema. Cells are filled
with placeholder values until the flagship run completes. The structure
--- which comparisons, which stratifications, which tests --- is fixed
ahead of time to avoid post-hoc selection.}

\subsection{In-domain perplexity}
\label{subsec:indomain-ppl}

\begin{table}[t]
\centering
\small
\begin{tabular}{lccccc}
\toprule
Model          & A1 & A2 & B1 & B2 & C1 \\
\midrule
B0 native GPT-2 & \textsc{tbd} & \textsc{tbd} & \textsc{tbd} & \textsc{tbd} & \textsc{tbd} \\
B1 learner GPT-2 & \textsc{tbd} & \textsc{tbd} & \textsc{tbd} & \textsc{tbd} & \textsc{tbd} \\
B2 + metadata concat & \textsc{tbd} & \textsc{tbd} & \textsc{tbd} & \textsc{tbd} & \textsc{tbd} \\
G1 ours         & \textsc{tbd} & \textsc{tbd} & \textsc{tbd} & \textsc{tbd} & \textsc{tbd} \\
\bottomrule
\end{tabular}
\caption{Perplexity on EFCAMDAT held-out, stratified by CEFR. The EFCAMDAT
test set contains A1--C1 only.
\textbf{Preliminary}: values pending the flagship run.}
\label{tab:ppl-cefr}
\end{table}

The predicted pattern: G1 should improve most on A1--A2 rows, where
baseline uncertainty is highest and metadata carries the most signal;
the gap should narrow on C1 rows where the native baseline is already
well-calibrated. B2 is expected to sit between B1 and G1.

\subsection{Cloze accuracy}
\label{subsec:cloze-results}

\begin{table}[t]
\centering
\small
\begin{tabular}{lcc}
\toprule
Model & Top-1 & Top-5 \\
\midrule
B0 & \textsc{tbd} & \textsc{tbd} \\
B1 & \textsc{tbd} & \textsc{tbd} \\
B2 & \textsc{tbd} & \textsc{tbd} \\
G1 & \textsc{tbd} & \textsc{tbd} \\
\bottomrule
\end{tabular}
\caption{Content-word cloze accuracy on held-out EFCAMDAT.
\textbf{Preliminary}: values pending the flagship run.}
\label{tab:cloze}
\end{table}

We predict that the gate helps cloze more than free-running NWP because
bidirectional context lets the metadata-conditioned filter exploit
proficiency signal more effectively: the A1 learner's distribution of
plausible completions for a masked determiner is narrower than the C1
learner's, and the gate can reshape the model's output distribution to
match.

\subsection{Zero-shot transfer}
\label{subsec:transfer}

\begin{table}[t]
\centering
\small
\begin{tabular}{lccc}
\toprule
Model & andrew100k & CELVA-SP & KUPA-KEYS \\
\midrule
B0 & \textsc{tbd} & \textsc{tbd} & \textsc{tbd} \\
B1 & \textsc{tbd} & \textsc{tbd} & \textsc{tbd} \\
B2 & \textsc{tbd} & \textsc{tbd} & \textsc{tbd} \\
G1 & \textsc{tbd} & \textsc{tbd} & \textsc{tbd} \\
\bottomrule
\end{tabular}
\caption{Zero-shot PPL on external corpora.
\textbf{Preliminary}: values pending the flagship run.}
\label{tab:transfer}
\end{table}

We expect the G1 gain to hold on andrew100k (same-target, different
register, fully-overlapping L1 inventory), hold on CELVA-SP (French is
in EFCAMDAT's inventory, so the gate's L1 embedding is usable), and
weaken on KUPA-KEYS (several L1s --- Polish, Greek, Hungarian, Dutch,
Bengali, Slovene --- are out of inventory and route to \texttt{unk}).
If the KUPA-KEYS gap is large, v07 (typological-feature gating)
becomes the path forward.
